\section{Generating $k$-Compositions of $n$ in Lexicographic Order}

Let us consider a non-maximal $k$-composition of $n$, $\gamma = (\gamma_1, \dots, \gamma_k)$. A composition $\gamma^{\prime}$ that is greater than $\gamma$ satisfies the property $\gamma^{\prime}_i > \gamma_i$ for the first index from the left $i$ where $\gamma$ and $\gamma^{\prime}$ differ. In order to minimize the increase from $\gamma$, we must find a $\gamma^{\prime}$ where $i$ is maximized. Furthermore, once we have maximized $i$, we need to minimize the part $\gamma^{\prime}_i$ subject to $\gamma^{\prime}_i > \gamma_i$. This is the foundation for generating the next lexicographic composition. It should be noted that the only information needed to generate such a composition is present within $\gamma$ itself.

Given $\gamma$, we wish to construct the next lexicographic composition $\gamma^{\prime}$. Whatever transformation we apply must preserve both the sum $n$ and the length $k$. We also know that we need to maximize the index $i$ where our next composition $\gamma^{\prime}$ differs from $\gamma$ and $\gamma^{\prime}_i > \gamma_i$. This implies we must focus our attention on the far right. In order to preserve the sum $n$, the index $i$ at which we increment must be accompanied by a decrement of the same magnitude at an index $j > i$.

Our first candidate index for incrementing is $k - 1$. Since we must decrement at an index greater than $k - 1$, we know that $\gamma_k$ must be decremented. Before applying these transformations, we must first check whether $\gamma_k > 1$. If $\gamma_k = 1$, decrementing would force $\gamma_k = 0$. Since we are not producing weak compositions, we would be violating the definition of only allowing parts to be members of $\mathbb{N}$. Therefore, if $\gamma_k = 1$, we know that $k - 1$ is an invalid index to increment.

Our next candidate index is $k - 2$. We only need to check if $\gamma_{k - 1} > 1$ as we have already determined $\gamma_{k} = 1$. If $\gamma_{k - 1} = 1$, we continue moving to the left by one. At each index $i$, we have already determined that the parts $(\gamma_{i + 2}, \dots, \gamma_k)$ are all equal to one in the previous iterations, so we only need to check $\gamma_{i + 1}$ to determine if we have reached a point where we can increment. Once we reach an index $i$ where $\gamma_{i + 1} > 1$, there are two cases we need to handle. The first case is when $i = k - 1$, which only requires that we increment $\gamma_{k - 1}$ and decrement $\gamma_k$. In the second case when $i < k - 1$, we increment $\gamma_i$, decrement $\gamma_{i + 1}$, and finally swap $\gamma_{i + 1}$ with $\gamma_k$. The last step is to ensure we produce the smallest lexicographic result. We know that after the decrement, $\gamma_{i + 1}$ is at least one and together with the fact that the remaining parts are all one, the swap guarantees that the possibly largest element is at the far right position.

\[
    \begin{aligned}
        \text{Before swap:} \quad & (\gamma_{i + 1} - 1, \dots, \gamma_k) = (\gamma_{i + 1} - 1, 1, \dots, 1) \\
        \text{After swap:} \quad & (1, \dots, 1, \gamma_{i + 1} - 1)
    \end{aligned}
\]

\begin{algorithm}
    \caption{Next Lexicographic $k$-Composition of $n$}

    \begin{algorithmic}[1] % [1] adds line numbers
        \Procedure{NextComposition}{$\gamma$}

            \Input a non-maximal $k$-composition $\gamma = (\gamma_1, \dots, \gamma_k)$ of $n$, where $k \geq 2$
            \Postcondition $\gamma$ is replaced by its next lexicographic successor

            \If{$\gamma_k > 1$}
                \State $\gamma_{k - 1} \gets \gamma_{k - 1} + 1$
                \State $\gamma_k \gets \gamma_k - 1$
            \Else
                \State $i \gets k - 2$

                \While{$i \geq 1 \land \gamma_{i + 1} = 1$}
                    \State $i \gets i - 1$
                \EndWhile

                \State $\gamma_i \gets \gamma_i + 1$
                \State $\gamma_{i + 1} \gets \gamma_{i + 1} - 1$
                \State $(\gamma_{i + 1}, \gamma_k) \gets (\gamma_k, \gamma_{i + 1})$
            \EndIf

        \EndProcedure
    \end{algorithmic}
\end{algorithm}

We will see in the next section that there are many challenges presented with the added restriction of distinctness. The key principles presented here provide the foundations for generating distinct compositions. There is an initial check for a common case. If this fails, we dive into a right to left scan where we try and find an index where the part isn't maximal. Finally, we apply a transformation to the tail that minimizes the difference from the current composition.
